
Language model evaluation typically produces aggregate accuracy metrics without actionable diagnostic signals. When a model fails on a benchmark question, the evaluation framework reports a binary outcome (correct/incorrect) but provides no indication of why the failure occurred or which intervention might remediate it. This diagnostic gap prevents targeted correction strategies and limits systematic improvement.

TruthfulQA \citep{Lin2021TruthfulQA} evaluates whether models produce truthful answers to questions designed to elicit common human falsehoods. State-of-the-art models achieve 60-70\% accuracy on this benchmark, yet the score alone reveals no information about failure modes. A model that incorrectly answers ''What is the capital of Australia?'' with ''Sydney'' has made an entity-substitution error --- confusing one Australian city for another --- which differs mechanistically from a reasoning failure or knowledge gap.

Existing evaluation pipelines aggregate failures into summary statistics, discarding the individual-level diagnostic information that could guide correction. Practitioners apply correction methods uniformly across all failures: retrieval-augmented generation (RAG) or chain-of-thought (COT) prompting are deployed without regard to whether the root cause involves factual retrieval, multi-step reasoning, or other failure modes.

This study addresses the integration gap between benchmark evaluation, interpretability analysis, and correction methods. We propose an automated classification framework that routes benchmark failures to attention-based diagnosis, enabling failure-type-specific correction routing without manual example selection or human annotation.

The framework operates as follows. First, named entity recognition identifies entity spans in benchmark questions. Second, attention weights over entity spans are extracted from the model's final transformer layer. Third, attention entropy is calculated as a measure of concentration versus diffusion. Fourth, a threshold classifier routes low-entropy failures to one correction method and high-entropy failures to another, based on the hypothesis that different attention patterns correspond to different failure modes.

Experiments on TruthfulQA single-entity factual questions using GPT-2 validate three key components:

\begin{enumerate}
\item Pre-validation conditions confirm that measurement assumptions hold (NER F1 = 0.96, Wikipedia coverage for entity-errors = 1.0).
\end{enumerate}

\begin{enumerate}
\item Entity-substitution errors exhibit significantly lower attention entropy (mean = 0.062) over NER-identified entity spans compared to non-entity errors (mean = 0.300), producing a robust statistical signal (p = 7.5e-07, Cohen's d = -1.13, N=73 processed samples).
\end{enumerate}

\begin{enumerate}
\item Threshold-based classification (threshold = 0.32) distinguishes entity-errors from non-entity errors with 86.7\% accuracy on held-out test data, exceeding the 70\% utility threshold and outperforming a random baseline by 33.4 percentage points.
\end{enumerate}

A synthetic correction experiment demonstrates matched routing structure using mock RAG and COT pipelines with configurable success rates. Results show +24 percentage points improvement for matched routing (entity-error to RAG) over mismatched routing (entity-error to COT) in GPT-3.5 mock tests, and +20 percentage points in Llama-2 mock tests. However, these experiments used synthetic data with pre-configured success rates and did not deploy real Wikipedia retrieval, LLM generation, or GPT-judge evaluation. Real-world correction effectiveness remains unvalidated.

Limitations include single-model validation (GPT-2 only; attention patterns in larger models unverified), manual gold labeling requirement (N=100 samples), 27\% sample loss due to tokenizer mismatch between spaCy NER and GPT-2 BPE, and single failure-type focus (entity-substitution only; reasoning errors not tested). The correction component validated pipeline structure but not deployed effectiveness.

The contributions are as follows:

\begin{itemize}
\item A method for automating attention-based failure diagnosis on 73 benchmark failures, exceeding the 10-20 manual example scale typical of prior interpretability work.
\end{itemize}

\begin{itemize}
\item Empirical validation of a robust attention entropy signature distinguishing entity-substitution errors from other failure modes (p < 0.001, large effect size).
\end{itemize}

\begin{itemize}
\item Demonstration that threshold-based classification enables automated failure routing with 86.7\% accuracy, providing a diagnostic signal without human annotation.
\end{itemize}

\begin{itemize}
\item Structural validation of matched correction routing framework using synthetic experiments; real-world deployment pending.
\end{itemize}

The study demonstrates that attention patterns provide actionable diagnostic signals at benchmark scale, enabling automated classification without per-instance human annotation. However, the correction effectiveness claim is not empirically validated and remains future work.
